\chapter{الفرائض المنصوصة وتفصيل السنة لمجمل القرآن}

\section{الفرائض المنصوصة التي سن رسول الله معها}
الحمد لله رب العالمين، والصلاة والسلام على رسول الله، أما بعد؛
يواصل الإمام الشافعي رحمه الله بيان موقع السنة من القرآن، ويذكر هنا أمثلة للفرائض التي نزل أصلها في القرآن، وجاءت السنة المطهرة لتوضحها وتبين كيفياتها وشروطها؛ إذ لو اكتفى الناس بتدبر القرآن لعرفوا الفرض، ولكن السنة أتت لتفصل وتبيّن.

\subsection{أحكام الوضوء والغسل (التفريق بين الفرض والسنة)}
قال الله تعالى: ﴿إِذَا قُمْتُمْ إِلَى الصَّلَاةِ فَاغْسِلُوا وُجُوهَكُمْ وَأَيْدِيَكُمْ إِلَى الْمَرَافِقِ وَامْسَحُوا بِرُءُوسِكُمْ وَأَرْجُلَكُمْ إِلَى الْكَعْبَيْنِ ۚ وَإِن كُنتُمْ جُنُبًا فَاطَّهَّرُوا﴾.
أبان القرآن أن طهارة الجنب تكون بالغسل، وأن طهارة الحدث الأصغر تكون بالوضوء. وسنّ رسول الله صلى الله عليه وسلم في غسل الجنابة أن يتوضأ المغتسل وضوء الصلاة أولاً، فهذا من سنته مع الغسل المنصوص.

وفي الوضوء، احتمل قوله تعالى: ﴿فَاغْسِلُوا وُجُوهَكُمْ﴾ غسلة واحدة أو أكثر، فجاءت السنة لتبيّن أن النبي صلى الله عليه وسلم توضأ (مرة مرة)، وتوضأ (مرتين مرتين)، وتوضأ (ثلاثاً ثلاثاً).
فدلت سنته حين اكتفى بالمرة الواحدة على أن الفرض الذي يجزئ في الوضوء هو المرة الواحدة، وأن ما زاد عليها (الثنتان والثلاث) هو نافلة وسنة مستحبة لطلب الفضل، وليس بفرض يفسد الوضوء بتركه.

كما بيّنت السنة دخول المرفقين والكعبين في الغسل، وأن الغسل يعمها؛ فقد كانت الآية تحتمل الغسل إليهما أو غسلهما، فبيّن فعله صلى الله عليه وسلم دخولهما.

\textbf{فائدة لغوية:} الوَضوء (بفتح الواو) هو الماء الذي يُتوضأ به، والوُضوء (بضم الواو) هو الفعل وحركة غسل الأعضاء.

\subsection{موانع الإرث (تخصيص السنة لعموم المواريث)}
قال تعالى: ﴿يُوصِيكُمُ اللَّهُ فِي أَوْلَادِكُمْ ۖ لِلذَّكَرِ مِثْلُ حَظِّ الْأُنثَيَيْنِ﴾، وقال في الكلالة (وهو من مات وليس له أصل ولا فرع وارث): ﴿وَإِن كَانَ رَجُلٌ يُورَثُ كَلَالَةً أَوِ امْرَأَةٌ وَلَهُ أَخٌ أَوْ أُخْتٌ فَلِكُلِّ وَاحِدٍ مِّنْهُمَا السُّدُسُ﴾.

هذه الآيات وردت بلفظ عام يشمل كل والد وولد وكل زوج وزوجة. فهل يرث القاتل عمداً? وهل يتوارث المسلم والكافر? وهل يرث العبد?
أتت السنة المطهرة لتخصص هذا العموم بضوابط وموانع، وبيّنت أن الميراث لا يثبت إلا بشروط:
\begin{itemize}
    \item \textbf{اتحاد الدين:} قال صلى الله عليه وسلم: «لا يَرِثُ الْمُسْلِمُ الْكَافِرَ وَلَا الْكَافِرُ الْمُسْلِمَ». فلا يتوارث أهل ملتين.
    \item \textbf{الحرية:} لأن العبد وما يملك هو ملك لسيده، فهو يباع ويوهب، فلا يرث ولا يُورث، ولو أورثناه لأعطينا المال لسيده الذي لا فرض له في كتاب الله.
    \item \textbf{البراءة من القتل العمد:} قال صلى الله عليه وسلم: «ليس لقاتل شيء»، فمن استعجل الشيء قبل أوانه عوقب بحرمانه، ويُمنع القاتل عمداً من الميراث عقوبة له.
\end{itemize}
فدلت السنة على أن الله تعالى إنما أراد بالآيات الخاصة بالميراث أناساً مخصوصين تتوفر فيهم هذه الشروط، دون غيرهم ممن شملهم اللفظ العام.

\subsection{تخصيص البيوع المحرمة من عموم (أحل الله البيع)}
قال تعالى: ﴿وَأَحَلَّ اللَّهُ الْبَيْعَ وَحَرَّمَ الرِّبَا﴾. هل كل البيوع حلال? 
لقد نهى رسول الله صلى الله عليه وسلم عن بيوع تراضى بها المتبايعان، ولم يكن فيها غرر مجهول، ولكنها تضمنت الربا المحرم. كأن يبيع الرجل الذهب القديم بذهب جديد ويدفع الفرق، فهذا بيع بالتراضي ولكنه محرم؛ لقوله صلى الله عليه وسلم: «الذهب بالذهب مثلاً بمثل ويداً بيد».

فدلت السنة الكريمة على أن الله تعالى إنما أحل البيوع التي لم يحرمها على لسان نبيه، فخصصت السنة عموم الآية، وبيّنت أن من البيوع ما هو باطل وإن تراضى عليه الناس.

\section{بيان السنة لمجمل الفرائض}
أحكم الله فرضه في كتابه مجملاً، كالصلاة والزكاة والحج، ثم بيّن على لسان نبيه صلى الله عليه وسلم كيفيات هذا الفرض:

\subsection{تفصيل أحكام الصلاة ومواقيتها}
قال تعالى: ﴿إِنَّ الصَّلَاةَ كَانَتْ عَلَى الْمُؤْمِنِينَ كِتَابًا مَّوْقُوتًا﴾. فأجمل الفرض. 
ثم بيّن النبي صلى الله عليه وسلم مواقيت الصلاة، وعددها بأنها خمس صلوات في اليوم والليلة لا يزاد عليها في الفرض. وبيّن أعداد ركعاتها، وأن الصبح ركعتان يجهر فيهما، والظهر والعصر أربع يسر فيهما، والمغرب ثلاث يجهر في اثنتين، والعشاء أربع يجهر في اثنتين. وبيّن أن الدخول فيها بالتكبير والخروج منها بالتسليم، وسنّ قصر الرباعية في السفر.

\subsection{صلاة الخسوف وصلاة الخوف}
وسنّ النبي صلى الله عليه وسلم في صلاة الكسوف ركعتين في كل ركعة ركوعين وسجودين، كما روى ذلك ابن عباس وعائشة رضي الله عنهما باختلاف يسير في الألفاظ واتفاق في الهيئة.

وكذلك في صلاة الخوف؛ فرغم أن الآية ذكرت صلاة الخوف مجملة: ﴿وَإِذَا كُنتَ فِيهِمْ فَأَقَمْتَ لَهُمُ الصَّلَاةَ فَلْتَقُمْ طَائِفَةٌ مِّنْهُم مَّعَكَ﴾، فقد صلاها النبي صلى الله عليه وسلم بهيئات متعددة تناسب حال الحرب وموقع العدو. ففي غزوة الخندق أخّر الصلاة حتى كُفي القتال فصلاها جمعاً، فلما نزلت صلاة الخوف نُسخ التأخير، وصار الواجب أداء الصلاة في وقتها بحسب الاستطاعة (رجالاً أو ركباناً، مستقبلي القبلة أو غير مستقبليها) في حال المسايفة واشتداد الخوف، ولا تُترك الصلاة أبداً لتُقضى بعد خروج وقتها إلا للنائم والناسي.

\section{خاتمة: مكانة السنة والرد على منكريها}
هذا الباب من أجلّ ما يُرد به على الطائفة المنحرفة التي تسمي نفسها بـ "القرآنيين"، وهم في الحقيقة مكذبون للقرآن؛ إذ لو اتبعوا القرآن حقاً لاتبعوا من أمره القرآن بطاعته. فالقرآن أتى بأصول الفرائض مجملة، ولولا سنة رسول الله صلى الله عليه وسلم ما عرف المسلمون كيف يصلون، ولا متى يزكون، ولا كيف يحجون، فمن ردّ السنة فقد هدم الدين كله وانسلخ منه.